\section{Certificate contract and independent checks}
\label{app:certificate}

The supplement constructs the $216\times216$ integer coboundary from the
cubical definition.  Its SHA-256 digest is
\begin{center}
\small\ttfamily
323c30c40d6d795ecf31aaaf3e22d681324449ff92a738a6f293c1ed2f72daee.
\end{center}
For every interior link it records selected plaquette indices, exact rational
covariance blocks, exact minimum-cost matrices, spectra, KKT ranks, and maximum
denominators of the witnesses.

The trust boundary is intentionally narrow.
\begin{enumerate}[label=(\arabic*)]
\item Floating singular-value and QR routines propose rational witnesses and
full-rank minors.
\item Rational reconstruction converts a projection lift to fractions.
\item Integer incidence multiplication verifies the image and orthogonality
conditions exactly.
\item Fraction-free elimination constructs the spill lifts.
\item The full rational KKT equations verify optimality.
\item Characteristic polynomials are expanded by the Leibniz formula over
\texttt{fractions.Fraction} and compared with the claimed factorization.
\end{enumerate}
If any proposal is wrong, an exact identity fails and no certificate is
emitted.  The final theorem does not depend on a tolerance used to label a
decimal eigenvalue.

One correction relative to the exploratory calculation deserves emphasis.
Centering the raw signed plaquette rows produces a different spectrum and a
spurious no-spill mode.  The certificate first multiplies each incident row by
the sign of the distinguished-link incidence.  Only then are the six rows
coordinates on one common staple transport space.  All results in this paper
use this convention-correct order $C_6O_\ell S_\ell$.
